\section{Trace bounds for ordered exponential products}

All norms in this section are operator norms on a finite-dimensional complex
Hilbert space. Products of matrices are ordered by increasing index.

\begin{theorem}[Multi-factor trace inequality]
    \label{thm:trace_holder_multi}
    \lean{Matrix.norm_trace_list_prod_mul_exp_smul_le}
    \leanok
    Let $B_0,\ldots,B_{m-1}$ be arbitrary operators, let
    $K_0,\ldots,K_{m-1}$ be Hermitian, and let $\theta_j\ge0$ satisfy
    $\sum_{j=0}^{m-1}\theta_j=1$. Then
    \begin{align}
        \left|\operatorname{tr}\prod_{j=0}^{m-1} B_j e^{\theta_jK_j}\right|
         & \le \prod_{j=0}^{m-1}\|B_j\|
        \prod_{j=0}^{m-1}\bigl(\operatorname{tr}e^{K_j}\bigr)^{\theta_j}.
        \label{eq:trace_holder_multi_bound}
    \end{align}
    A zero weight contributes an identity exponential. On the zero-dimensional
    space the assertion reads $0\le0$, with $0^0=1$ in the scalar factors.
\end{theorem}
\begin{proof}\leanok
    Prove the stronger statement with complex weights $w_j$, nonnegative real
    parts summing to one, and the same right-hand side with exponents
    $\operatorname{Re}w_j$. The imaginary exponential $e^{i\operatorname{Im}(w_j)K_j}$
    is unitary and can be absorbed into $B_j$ without changing its norm.
    Induct on the number of factors. If the remaining real weights sum to
    zero, cyclicity of the trace reduces the estimate to
    $|\operatorname{tr}(PC)|\le\operatorname{tr}(P)\|C\|$ for $P\succeq0$.

    Otherwise put $S=1-\operatorname{Re}w_0>0$. In the strip
    $0\le\operatorname{Re}z\le1$, replace the first weight and, for $1\le j<m$, the remaining weights by
    \begin{align}
        w_0(z) & =z+i\operatorname{Im}w_0,
               & w_j(z)                    & =\frac{1-z}{S}\operatorname{Re}w_j+i\operatorname{Im}w_j
        .
        \label{eq:trace_holder_multi_interpolation}
    \end{align}
    On the boundary $\operatorname{Re}z=0$, the first exponential is unitary;
    merge it with the next factor and apply induction. On the boundary
    $\operatorname{Re}z=1$, use the preceding positive-matrix estimate.
    The trace of the product is holomorphic. Its modulus is bounded on the
    closed strip: remove the unitary factors and bound the remaining norms
    continuously in $\operatorname{Re}z\in[0,1]$.
    Hadamard's three-lines theorem at $z=\operatorname{Re}w_0$ now gives
    (\ref{eq:trace_holder_multi_bound}), since the powers of each operator
    norm add to one and the trace powers become $\operatorname{Re}w_j$.
\end{proof}

\begin{theorem}[Imaginary-time trace inequality]
    \label{thm:trace_holder_imaginary_time}
    \lean{Matrix.IsHermitian.norm_trace_list_prod_mul_exp_neg_smul_le}
    \leanok
    Let $H$ be Hermitian, let $B_0,\ldots,B_k$ be arbitrary operators,
    and let $a_j\ge0$ with $\sum_{j=0}^k a_j=\beta$. Then
    \begin{align}
        \left|\operatorname{tr}\prod_{j=0}^k B_j e^{-a_jH}\right|
         & \le \left(\prod_{j=0}^k\|B_j\|\right)
        \operatorname{tr}e^{-\beta H}.
        \label{eq:trace_holder_imaginary_time}
    \end{align}
    This is the unnormalized trace bound used
    in~\cite{Hastings2006Solving}, arXiv:cond-mat/0508554v3,
    source lines 1141--1151. For $k\ge1$ and ordered times
    $0\le\tau_1\le\cdots\le\tau_k\le\beta$, take successive differences
    $a_0=\tau_1$, $a_j=\tau_{j+1}-\tau_j$ for $1\le j<k$, and $a_k=\beta-\tau_k$.
    On a nonzero-dimensional space, division by the positive partition
    function gives the source's normalized inequality. No commutation of
    the $B_j$ with $H$ is required.
\end{theorem}
\begin{proof}\leanok
    \uses{thm:trace_holder_multi}
    Nonnegativity of the $a_j$ gives $\beta\ge0$. If $\beta>0$, apply
    Theorem~\ref{thm:trace_holder_multi} with $K_j=-\beta H$ and
    $\theta_j=a_j/\beta$. The trace powers combine to one.
    If $\beta=0$, all $a_j$ vanish; apply the same theorem with $K_j=0$
    and weight one on the first factor, zero on the others.
\end{proof}
